% Generated draft: requires mathematical and source review.
\documentclass[11pt]{article}
\usepackage[margin=25mm]{geometry}
\usepackage{fontspec}
\setmainfont{Latin Modern Roman}
\usepackage{amsmath,amssymb,mathtools,unicode-math}
\setmathfont{Latin Modern Math}
\usepackage{enumitem,needspace,xurl,hyperref}
\setlength{\parindent}{0pt}
\setlength{\parskip}{5pt}
\setlength{\emergencystretch}{3em}
\newcommand{\N}{\mathbb{N}}
\newcommand{\Z}{\mathbb{Z}}
\newcommand{\Q}{\mathbb{Q}}
\newcommand{\R}{\mathbb{R}}
\newcommand{\C}{\mathbb{C}}
\newcommand{\F}{\mathbb{F}}
\newcommand{\A}{\mathbb{A}}
\newcommand{\PP}{\mathbb{P}}
\newcommand{\E}{\mathbb{E}}
\newcommand{\eps}{\varepsilon}
\newcommand{\dd}{\,\mathrm{d}}
\begin{document}
\section*{UnsolvedMath ID: 6200068}
\textbf{Source problem code:} catalogue entry
\subsection{Definitions and mathematical statement}
Let $F$ be a free group and $\mu$ a probability measure on it. For the right random walk started at the identity, let $X_0=e$ and $X_n=Y_1\cdots Y_n$, where the increments $Y_i$ are independent with law $\mu$. Its Markov operator is
\[
(P_\mu u)(g)=\sum_{h\in F}\mu(h)u(gh),
\]
and a bounded function is $\mu$-harmonic when $P_\mu u=u$.

Determine the Poisson boundary for arbitrary $\mu$. One probabilistic definition uses the random-walk path space with its Markov law and the sigma-algebra of events invariant under deleting the initial state, modulo null events. The corresponding measure-space factor is the Poisson boundary; it encodes asymptotic boundary data for bounded harmonic functions. For a non-generating measure, the reachable states and possible distinct communicating classes must be accounted for rather than assuming irreducibility.

The question imposes no finite-support, symmetry, generating-support, entropy, or moment hypothesis on $\mu$. The rank of the free group is not specified either. A solution should therefore state exactly which free groups and measures its boundary model covers, including degenerate support cases. The supplied summary cites known special cases but gives no hypotheses precise enough to turn its references to ends, horospheric models, or Gromov boundaries into a classification for all measures.
\subsection{Short English statement}
Describe the Poisson boundary of a random walk on a free group for every probability measure, without silently adding support, symmetry, entropy, or moment assumptions.
\subsection{Sources}
\begin{itemize}
\item Misha Kapovich, \emph{Problems on Boundaries of Groups and Kleinian Groups} (2005), Problem 68, PDF page 18, attributed to Vadim Kaimanovich, \url{https://www.math.ucdavis.edu/~kapovich/EPR/problems.pdf}; accessed 2026-07-29, PDF-text-normalized extraction.
\item Related-work names in the supplied research summary: Ledrappier, Kaimanovich, Kaimanovich--Vershik, Derriennic, Gaboriau, and Willis. The record gives no publication identifiers or complete hypotheses for those attributions.
\item Local registry \texttt{data/unsolvedmath/problems.json}, permanent UnsolvedMath ID $6200068$, source code catalogue entry, \emph{Boundaries of Groups and Kleinian Groups --- Problem 68}. Dataset attribution: Ulam AI and the UnsolvedMath Contributors.
\end{itemize}
\end{document}
